\section{Runtime Interfaces and Data Contracts}
\label{sec:runtime-contracts}
% ============================================================

The implemented components exchange structured messages so that a language-model output cannot become a robot action without validation. Each contract below identifies the data that crosses one implementation boundary and the component that produces, checks, or consumes it.

\subsection{Runtime Contract Overview}

Table~\ref{tab:contract-map} summarises each runtime contract, its interface or payload, and its owner.

\begin{table}[H]
\centering
\scriptsize
\renewcommand{\arraystretch}{1.18}
\caption{Runtime contract map and ownership in the active dialogue, planning, and execution stack.}
\label{tab:contract-map}

\begin{tabularx}{\textwidth}{
@{}
>{\raggedright\arraybackslash}p{0.04\textwidth}
>{\raggedright\arraybackslash}p{0.17\textwidth}
>{\raggedright\arraybackslash}p{0.38\textwidth}
>{\raggedright\arraybackslash}X
@{}}
\toprule
\textbf{ID} & \textbf{Contract} & \textbf{Interface / payload} & \textbf{Ownership}\\
\midrule

1 & Chatbot turn output
& Internal chatbot result before ROS planner publication.
& \code{dialogue_manager} invokes the chatbot path; \code{chatbot_llm} produces the semantic result.\\

2 & Planner request and admission
& Candidate topic \code{/nao_orchestrator/planner_request}. Admitted topic \code{/planner/request}. Planner-facing payload is carried in \code{Intent.data}.
& \code{chatbot_llm} proposes the execution request; \code{nao_orchestrator} admits and normalizes it; \code{planner_llm} consumes the admitted request.\\

3 & Compact grounding
& \code{grounded_context} inside Contract~2.
& Scene and KB projection provide compact symbolic context for \code{chatbot_llm} and \code{planner_llm}.\\

4 & Shared skill entry
& Registry projection consumed by the planner and orchestrator; its adapter mapping identifies the receiving skill server.
& \code{skill_common} defines the shared capability shape; \code{planner_llm} uses it for prompt construction and \code{nao_orchestrator} uses it for plan validation and dispatch.\\

5 & Planner output
& Topic \code{/intents}. Structured plan payload in \code{Intent.data.plan}.
& \code{planner_llm} produces the plan; \code{nao_orchestrator} validates and dispatches executable steps.\\

6 & Skill execution result
& ROS action result or service response with a typed skill-result payload.
& \code{nao_orchestrator} dispatches the call; the skill server owns fresh evidence; verified world-state updates are applied through \code{kb_skills}.\\

7 & Execution feedback
& Topic \code{/planner/execution_feedback}. Structured JSON feedback payload.
& \code{nao_orchestrator} publishes feedback; a goal-keyed planning-state helper inside \code{planner_llm} uses it for replanning, clarification, or failure handling.\\

8 & Planner dialogue act
& Topic \code{/planner/dialogue_act}. Structured communication request.
& \code{planner_llm} emits dialogue-act intent; \code{nao_orchestrator} relays it; \code{dialogue_manager} realizes the final speech.\\

9 & RDF knowledge state and raw evidence
& KnowledgeCore statements through \code{/kb/query} and \code{/kb/revise}; source evidence includes tracked-person state, \code{/scene/summary}, and skill results.
& Source components own fact freshness; \code{kb_skills} owns KB transport; compact projections and feedback preserve provenance without replacing raw evidence.\\

\bottomrule
\end{tabularx}
\end{table}

The overview keeps only the fields needed to identify each contract. Later sections add the runtime metadata needed for orchestration, feedback-driven planning, tracing, and debugging.

\subsection{Contract 1: Chatbot Turn Output}

The intent in this contract is the user's goal as extracted by the chatbot. It is not a planner output.

\begin{lstlisting}
{
  "verbal_ack": "Okay, I will do that.",
  "route": "execution",
  "confidence": 0.82,
  "user_intent": {
    "type": "look_at",
    "goal_text": "look at the person",
    "scene_targets": ["person"],
    "request_kind": "new_goal"
  }
}
\end{lstlisting}

Table~\ref{tab:chatbot-fields} identifies the fields that remain visible after the chatbot completes one semantic routing decision.

\begin{table}[H]
\centering
\footnotesize
\renewcommand{\arraystretch}{1.18}
\caption{Fields in Contract 1: chatbot turn output.}
\label{tab:chatbot-fields}

\begin{tabularx}{\textwidth}{
@{}
>{\raggedright\arraybackslash}p{0.33\textwidth}
>{\raggedright\arraybackslash}X
@{}}
\toprule
\textbf{Field} & \textbf{Meaning}\\
\midrule

\code{verbal_ack}
& Optional immediate acknowledgement generated by the chatbot layer. This may be spoken or logged before execution begins, depending on the dialogue policy.\\

\code{route}
& Turn classification produced by \code{chatbot_llm}. Supported values include \code{dialogue}, \code{knowledge_query}, and \code{execution}. Only \code{execution} turns enter the planner path; \code{dialogue} and \code{knowledge_query} remain non-execution turns but are still useful for tracing.\\

\code{confidence}
& Confidence assigned by the chatbot or intent-extraction stage. It is useful for debugging and possible admission policy, but it is not a planner priority score.\\

\code{user_intent.type}
& Normalized user-intent label extracted by \code{chatbot_llm}, such as \code{look_at}, \code{find_object}, \code{scan}, or \code{navigate_to}.\\

\code{user_intent.goal_text}
& Natural-language objective passed toward the planner path as the task to satisfy. This field preserves the user's goal in a compact form.\\

\begin{tabular}[t]{@{}l@{}}
\code{user_intent.}\\
\code{scene_targets}
\end{tabular}
& Optional target hints extracted from the utterance, such as \code{person}, \code{cup}, or \code{table}. These are hints, not validated grounded entities.\\

\begin{tabular}[t]{@{}l@{}}
\code{user_intent.}\\
\code{request_kind}
\end{tabular}
& Proposed goal transition class, such as \code{new_goal}, \code{update_goal}, \code{answer_to_clarification}, or \code{cancel_goal}. The planner gate inside \code{nao_orchestrator} should normalize and admit this field before publishing \code{/planner/request}.\\

\bottomrule
\end{tabularx}
\end{table}


\subsection{Contract 2: Planner Request}

The planner request is the task description generated by the orchestrator from the chatbot output and consumed by the planner.

\begin{lstlisting}
{
  "goal_id": "goal_turn_123",
  "request_kind": "new_goal",
  "goal_text": "look at the person and report completion",
  "normalized_intents": ["look_at"],
  "grounded_context": {
    "entities": [
      {
        "id": "anonymous_person_ehfbf",
        "label": null,
        "kind": "person",
        "class": "Human",
        "visible": true,
        "relations": []
      }
    ]
  }
}
\end{lstlisting}

Table~\ref{tab:planner-request-fields} separates request identity, goal semantics, and grounded evidence before the planner is invoked.

\begin{table}[H]
\centering
\footnotesize
\renewcommand{\arraystretch}{1.18}
\caption{Fields in Contract 2: minimal planner-facing request.}
\label{tab:planner-request-fields}

\begin{tabularx}{\textwidth}{
@{}
>{\raggedright\arraybackslash}p{0.33\textwidth}
>{\raggedright\arraybackslash}X
@{}}
\toprule
\textbf{Field} & \textbf{Meaning}\\
\midrule

\code{goal_id}
& Identifier for the active objective. It groups replans and feedback for the same task.\\

\code{request_kind}
& States whether this is a new goal, an update to the current goal, a cancellation, or another supported transition. \code{chatbot_llm} may propose it, but \code{nao_orchestrator} normalizes and admits it before publishing \code{/planner/request}.\\

\code{goal_text}
& Objective extracted from the user utterance. This is the central planning input for \code{planner_llm}.\\

\code{normalized_intents}
& Intent labels extracted by \code{chatbot_llm}. They help \code{planner_llm} choose suitable skills from the available skill registry.\\

\code{grounded_context}
& Compact scene graph used by the planner as the current symbolic state of the world. It contains validated entities and relations rather than raw detector output.\\

\bottomrule
\end{tabularx}
\end{table}

\subsection{Contract 3: Grounded Context}

The planner receives a compact, JSON-native view of relevant world state. This is best described as a compact scene graph derived from the KB, continually updated by the perception pipeline (human/person detection and object perception). Raw KB triples, detector metadata, and tracker-derived details remain available through scene/KB evidence paths to skills, but they are not duplicated into every LLM prompt.

\Needspace{18\baselineskip}
\begin{lstlisting}
{
  "grounded_context": {
    "entities": [
      {
        "id": "cup_jrjic",
        "label": "cup",
        "kind": "object",
        "class": "Cup",
        "visible": true,
        "relations": [
          {"predicate": "dbp:color", "object": "blue"},
          {"predicate": "oro:isOn", "object": "table_1"}
        ]
      },
      {
        "id": "anonymous_person_ehfbf",
        "label": null,
        "kind": "person",
        "class": "Human",
        "visible": true,
        "relations": []
      }
    ]
  }
}
\end{lstlisting}

Table~\ref{tab:grounded-context-fields} defines the compact entity and relation fields used by dialogue and planning.

\begin{table}[H]
\centering
\footnotesize
\renewcommand{\arraystretch}{1.18}
\caption{Fields in Contract 3: grounded context.}
\label{tab:grounded-context-fields}

\begin{tabularx}{\textwidth}{
@{}
>{\raggedright\arraybackslash}p{0.33\textwidth}
>{\raggedright\arraybackslash}X
@{}}
\toprule
\textbf{Field} & \textbf{Meaning}\\
\midrule

\code{id}
& Stable entity handle used by planner steps and skill arguments. This is the preferred reference when a skill needs to target a known person, object, or location.\\

\code{label}
& Human-readable name or class hint. It may be empty for anonymous people or entities that should not be assigned a personal name.\\

\code{kind}
& Entity category, such as \code{object}, \code{person}, or another supported runtime category. This prevents people from being treated as generic objects.\\

\code{class}
& Compact semantic class used by the planner for reasoning, such as \code{Cup}, \code{Human}, or \code{Table}.\\

\code{visible}
& Indicates whether the entity is currently grounded in the latest accepted scene or KB state. It should not be treated as a guarantee that the entity will still be available at execution time.\\

\code{relations}
& Bounded list of relevant semantic relations, such as colour, containment, support, proximity, or ownership-style facts. These relations add task-relevant context without exposing the full KB.\\

\bottomrule
\end{tabularx}
\end{table}

\subsection{Contract 4: Shared Skill Entry}

The planner and orchestrator use one shared entry shape to describe a skill before it is placed in a prompt, admitted for execution, or mapped to an action server. The entry separates the information needed to select a skill from the information needed to verify its result.

\begin{lstlisting}
{
  "object_id": "navigate_to",
  "kind": "skill",
  "category": "navigation",
  "aliases": ["go_to", "move_to_location"],
  "params": ["target", "location"],
  "required_params": ["target"],
  "preconditions": ["navigation is enabled",
                    "target is known or resolvable"],
  "expected_effects": ["robot arrives at the target location"],
  "observable_success": ["status"],
  "failure_modes": ["path_blocked", "unknown_location",
                    "safety_disabled"],
  "planner_guidance": ["use for a semantic destination"],
  "result_schema": {"required": ["skill", "status", "summary_text"]},
  "safety_flags": ["navigation"],
  "implementation_status": "fake",
  "robot_adapter_mapping": "fake_skills.navigate_to",
  "supports_failure_injection": true,
  "metadata": {"runtime_callable": true}
}
\end{lstlisting}

Table~\ref{tab:skill-entry-contract} names the common capability fields consumed by prompt construction, admission, dispatch, and result interpretation.

\begin{table}[H]
\centering
\footnotesize
\renewcommand{\arraystretch}{1.15}
\caption{Shared skill-entry fields used by planning and execution.}
\label{tab:skill-entry-contract}
\begin{tabularx}{\textwidth}{
@{}
>{\raggedright\arraybackslash}p{0.25\textwidth}
>{\raggedright\arraybackslash}X
@{}}
\toprule
\textbf{Field} & \textbf{Runtime role}\\
\midrule
\code{object_id}, \code{kind}, \code{category}, \code{aliases}
& Canonical planner name, semantic group, and accepted compatibility names.\\
\code{params}, \code{required_params}
& Typed planner-facing arguments and the subset required before admission.\\
\code{preconditions}
& Grounding or state assumptions that must be satisfied or explicitly owned by a scenario. They do not prove a future effect.\\
\code{expected_effects}
& Intended state changes or evidence products used to select and order skills.\\
\code{observable_success}
& Result or trace evidence that can support completion after execution.\\
\code{failure_modes}
& Typed outcomes used by the orchestrator and the goal-state helper to stop, clarify, or request replanning.\\
\code{planner_guidance}
& Bounded selection and ordering guidance inserted into the planner prompt.\\
\code{result_schema}
& Required and optional result fields, including the summary and evidence payload.\\
\code{safety_flags}
& Safety-relevant labels used during admission and dispatch.\\
\path{metadata.runtime_callable}
& Whether the registry permits planner-visible runtime use.\\
\code{implementation_status}
& Current implementation class, such as first-party, fake, or hybrid.\\
\code{robot_adapter_mapping}
& Adapter or skill server that receives the dispatched call.\\
\code{supports_failure_}\newline\code{injection}
& Whether controlled failure outcomes are available for validation.\\
\bottomrule
\end{tabularx}
\end{table}

The same shape is projected to the planner prompt, deterministic orchestrator checks, controlled fake skill simulation policies, and adapter configuration. Preconditions constrain admission, while expected effects remain hypotheses until a skill returns observable evidence. This distinction is used by \code{report_result}: a final summary must be derived from result payloads or trace evidence rather than copied from a planner-generated claim. The complete runtime inventory is provided in Appendix~\ref{sec:runtime-skill-inventory}.

\subsection{Contract 5: Planner Output}

The planner output is a plan: a sequence of skill steps to achieve the user goal. It should focus on executable content and minimal lineage. Communication policy and retry metadata are still valid runtime fields when used by the stack, but this controlled version foregrounds the ordered \code{steps} list because that is the essential planner output.

\Needspace{18\baselineskip}
\begin{lstlisting}
{
  "plan": {
    "goal_id": "goal_turn_123",
    "plan_id": "plan_123",
    "plan_version": 1,
    "status": "planned",
    "steps": [
      {
        "id": "step_1",
        "type": "skill",
        "name": "look_at",
        "args": {"target_frame": "anonymous_person_ehfbf"},
        "requires": [],
        "on_failure": "replan"
      }
    ]
  }
}
\end{lstlisting}

Table~\ref{tab:planner-output-fields} identifies the plan and step fields that must survive deterministic validation before dispatch.

\begin{table}[H]
\centering
\footnotesize
\renewcommand{\arraystretch}{1.18}
\caption{Fields in Contract 5: planner output.}
\label{tab:planner-output-fields}

\begin{tabularx}{\textwidth}{
@{}
>{\raggedright\arraybackslash}p{0.32\textwidth}
>{\raggedright\arraybackslash}X
@{}}
\toprule
\textbf{Field} & \textbf{Meaning}\\
\midrule

\code{goal_id}
& User objective that this plan attempts to satisfy. It links the planner output to the active goal.\\

\code{plan_id}
& Identifier for this concrete plan proposal. It is useful when the runtime keeps plan history or compares revised plans for the same goal.\\

\code{plan_version}
& Revision number for replans of the same goal. This helps distinguish the initial plan from later feedback-driven planning outputs.\\

\code{status}
& Planning outcome for this plan, such as \code{planned}, \code{needs_clarification}, or \code{failed_to_plan}. This describes whether \code{planner_llm} produced an executable plan, not whether the robot has already executed it.\\

\code{steps}
& Ordered list of skill calls produced by the planner. This is the essential planner output consumed by \code{nao_orchestrator}.\\

\code{steps[].id}
& Stable step identifier used for execution feedback, dependency tracking, and possible replan joins.\\

\code{steps[].name}
& Skill name selected from the allowed skill registry. The planner should not invent skills or use robot-specific APIs directly.\\

\code{steps[].args}
& Skill arguments. They should refer to known entity IDs, explicit primitive values, or arguments supported by the selected skill schema.\\

\code{steps[].requires}
& Dependencies between steps, if any. This allows the runtime to know which steps depend on previous successful outcomes.\\

\begin{tabular}[t]{@{}l@{}}
\code{steps[].}\\
\code{on_failure}
\end{tabular}
& Recovery policy proposed by \code{planner_llm} for this step. If the step fails, \code{nao_orchestrator} may fail safely, continue if allowed, request user input, or send feedback back to \code{planner_llm} for replanning. Allowed values include \code{replan}, \code{clarify}, \code{ask_user}, \code{continue}, \code{ignore}, and \code{fail}.\\

\bottomrule
\end{tabularx}
\end{table}

\paragraph{Runtime recovery metadata.}
The current contract supports recovery metadata as skill failure may require controlled recovery. Each step may carry a \code{retry_budget}, and the plan envelope may also carry a plan-level \code{retry_budget}. This does not mean that \code{retry} is a separate \code{on_failure} value. Instead, retry behaviour is controlled by \code{nao_orchestrator} using the available retry budget together with the step's \code{on_failure} policy. The plan envelope may also carry \code{communication_policy}. This controls when the planner/dialogue path should communicate progress, completion, clarification needs, or failure.

When a step fails with \code{on_failure = replan}, \code{nao_orchestrator} halts or blocks continuation, publishes execution feedback to \code{planner_llm}, and waits for a revised validated plan before dispatching further actions. This keeps execution authority inside \code{nao_orchestrator}, while allowing recovery to be driven by feedback-driven planning rather than by hard-coded local replanning.

\subsection{Contract 6: Skill Execution Result}

Every dispatched skill returns a typed result before the orchestrator can mark its step as completed. The result distinguishes the execution status from the evidence that supports it. For state-changing skills, the evidence may also describe world-state changes that must be reflected in the KB. These changes are reported by the skill that observed or produced them; they are not inferred from the planner's expected-effects text in Contract~4.

\begin{lstlisting}
{
  "skill": "place_object",
  "status": "succeeded",
  "summary_text": "I placed cup_1 on table_1.",
  "evidence": {
    "object_id": "cup_1",
    "destination_id": "table_1",
    "kb_effects": [
      {"operation": "remove", "statement": "robot oro:holds cup_1"},
      {"operation": "add", "statement": "cup_1 oro:isOn table_1"}
    ]
  },
  "failure": null,
  "metadata": {"fake": true}
}
\end{lstlisting}

The implementation retains the field name \code{evidence.kb_effects}, but the field describes observed world-state changes that must be reflected in the KB and the orchestrator applies these statements only after a successful result.

\subsection{Contract 7: Execution Feedback}

Execution feedback tells the planner which goal and step were affected, whether the step succeeded or failed, and why. This compact information is sufficient to decide whether to continue, replan, request clarification, or report failure.

\Needspace{18\baselineskip}
\begin{lstlisting}
{
  "goal_id": "goal_turn_123",
  "status": "failed",
  "step": {
    "id": "step_1",
    "name": "look_at"
  },
  "reason": "target frame unavailable",
  "result_payload": {
    "target_found": false,
    "target_kind": "person"
  }
}
\end{lstlisting}



Runtime feedback may include additional fields used by \code{nao_orchestrator}, logging, and the goal-state helper. Lineage uses \code{plan_id} and \code{plan_version}; event handling may add \code{event_type}, \code{validation_status}, \code{blocking}, and \path{needs_user_input}; observability may add \code{timestamp_sec}, \code{result_summary}, and retry metadata. These fields support execution control and debugging, while the minimal planner-facing feedback remains the goal, step, status, reason, and result payload.

\subsection{Contract 8: Planner Dialogue Act}

Planner dialogue acts represent planner intent to communicate, not direct speech execution. They are relayed through the orchestrator and realized by the dialogue layer.

\begin{lstlisting}
{
  "goal_id": "goal_turn_123",
  "act": "ask_clarification",
  "await_user_response": true,
  "reason": "multiple visible people",
  "text_hint": "Which person should I look at?",
  "slots_needed": ["target_person"]
}
\end{lstlisting}

Allowed dialogue acts are acknowledgement, progress update, clarification, help request, failure explanation, completion notification, and cancellation notification. The planner does not speak directly; it emits a structured communication request.

\subsection{Contract 9: RDF Knowledge State}

The implementation reuses the open-source KnowledgeCore package rather than introducing a new knowledge-base system \parencite{knowledgecore}. KnowledgeCore represents symbolic world state as an RDF-style graph. Each statement has a subject, predicate, and object. Stable subjects identify people, objects, places, supports, and the robot; predicates express class membership, attributes, visibility, containment, and spatial relations. The same statement shape is used by preloaded fixtures, detector-derived object facts, explicit user-authorised KB mutations, and verified post-skill effects.

\begin{lstlisting}
codex_lab_cup rdf:type Cup
codex_lab_cup dbp:color red
codex_lab_cup oro:isOn codex_lab_table
codex_lab_room oro:contains codex_lab_table
person_alex rdf:type Human
myself sees person_alex
\end{lstlisting}

The principal predicate families are shown in Table~\ref{tab:rdf-kb-predicates}. A class statement such as \code{rdf:type Human} also preserves the semantic distinction between a person and an object even when both can be action targets.

\begin{table}[H]
\centering
\footnotesize
\renewcommand{\arraystretch}{1.15}
\caption{Principal RDF-style predicate families used by the runtime KB.}
\label{tab:rdf-kb-predicates}
\begin{tabularx}{\textwidth}{@{}p{0.24\textwidth}p{0.29\textwidth}X@{}}
\toprule
\textbf{Role} & \textbf{Representative predicates} & \textbf{Runtime meaning}\\
\midrule
Entity class & \code{rdf:type} & Distinguishes people, objects, supports, rooms, and other planner-relevant classes.\\
Attributes & \code{dbp:name}, \code{dbp:color} & Provides stable semantic labels and disambiguating properties.\\
Observation & \code{sees}, visual-centre, score, source, and last-seen predicates & Records who observed an entity and the bounded provenance/freshness of a detector observation.\\
Spatial relations & \code{oro:isOn}, \code{oro:isAt}, \code{oro:isIn}, \code{oro:contains} & Represents support, destination, room membership, and containment used by grounding and skill effects.\\
\bottomrule
\end{tabularx}
\end{table}

KB mutations use add, revise, and remove operations on this shared statement form through \code{/kb/revise}; reads use query patterns through \code{/kb/query}. Transient perception facts carry a finite lifespan, while explicit fixture and verified skill-effect facts follow the lifetime declared by their owner. Reused person-perception components publish person identities and tracked state through their public interfaces \parencite{mohamed2020ros4hri}. Object detections are handled separately by \code{nao_scene_grounding}, which publishes \code{/scene/summary} and detector-derived KB statements.

% ============================================================
